% Written from cas-sc-template.tex of the Elsevier CAS bundle 2.4 (class cas-sc, single column).

% float (for the [H] figures and the [H] algorithm) is loaded before the class: the class loads hyperref,
% which adapts its float anchors only to a float package loaded before it (otherwise each figure anchor is
% written twice). \usepackage{float} in the package list below is then a no-op.
\RequirePackage{float}
\documentclass[a4paper,fleqn]{cas-sc}

\usepackage[numbers,sort&compress]{natbib}

% Packages the text needs beyond those the class loads.
\usepackage{float}            % H placement of figures and algorithms
\usepackage[most]{tcolorbox}
\usepackage{amsthm}
\usepackage{algorithm}
\usepackage{algpseudocode}    % part of the algorithmicx package

% Theorem-like environments, the "problem" box and numbering; independent of the document class.

% Numbering within sections
\numberwithin{equation}{section}

\theoremstyle{definition}
\newtheorem{definition}{Definition}[section]

\theoremstyle{plain}
\newtheorem{theorem}[definition]{Theorem}
\newtheorem{lemma}[definition]{Lemma}
\newtheorem{claim}[definition]{Claim}
\newtheorem{corollary}[definition]{Corollary}
\newtheorem{proposition}[definition]{Proposition}

\theoremstyle{remark}
\newtheorem{remark}[definition]{Remark}
\newtheorem{example}[definition]{Example}

% "problem" environment: framed statement of a computational problem
\newtcolorbox{problem}[1][]{%
  enhanced,
  colback=white,      % background color
  colframe=black,     % border color
  boxrule=0.8pt,      % border thickness
  arc=2pt,            % corner roundness
  left=6pt, right=6pt, top=6pt, bottom=6pt,
  title=\textbf{Problem: #1},
}

% Algorithm headers: every algorithm states its Input and Output;
% algpseudocode would print "Require:" and "Ensure:".
\algrenewcommand\algorithmicrequire{\textbf{Input:}}
\algrenewcommand\algorithmicensure{\textbf{Output:}}

% Notation macros; independent of the document class.

% Notation
\newcommand{\cp}{\operatorname{cp}}                  % value of a minimum cut-path
\newcommand{\CP}{\operatorname{CP}}                  % set of all cut-paths
\newcommand{\diam}{\operatorname{diam}}
\newcommand{\OPT}{\mathrm{OPT}}
\newcommand{\MinCutPath}{\textsc{Min Cut-Path}}
\newcommand{\SSP}{\textsc{Separating Shortest Path}}
\newcommand{\ThreeSAT}{\textsc{3-SAT}}

% The class sets the text in STIX, which has no italic small caps: \textsc inside italic text takes the
% upright small caps that the font selection substitutes anyway; declared here so that it gives no warning.
\AtBeginDocument{%
  \DeclareFontShape{T1}{stix}{m}{scit}{<->ssub*stix/m/sc}{}%
  \DeclareFontShape{T1}{stix}{b}{scit}{<->ssub*stix/b/sc}{}}

% Key nologo of the class: the e-mail address is labelled "Email address:" instead of the icon
% thumbnails/cas-email.jpeg, so the manuscript also builds from one folder without sub-folders.
% Key H: the class's figure takes key-value options and would lose a plain [H] (it stores an empty placement);
% H is declared as a short form of the class's own key pos=H, so the figures keep the standard \begin{figure}[H].
\ExplSyntaxOn
\keys_set:nn { stm / mktitle } { nologo }
\keys_define:nn { cas / fig } { H .meta:n = { pos = H } }
\ExplSyntaxOff

\begin{document}
\let\WriteBookmarks\relax
\def\floatpagepagefraction{1}
\def\textpagefraction{.001}

\shorttitle{Min Cut-Path Problem}

\shortauthors{N. Micheľ}

\title[mode = title]{Min Cut-Path Problem}

\author[1]{Norbert Micheľ}

\cormark[1]

\ead{5344553@upjs.sk}

\affiliation[1]{organization={Institute of Computer Science, Faculty of Science, Pavol Jozef Šafárik University in Košice},
            addressline={Jesenná 5},
            postcode={040 01},
            postcodesep={},
            city={Košice},
            country={Slovakia}}

\cortext[1]{Corresponding author}

% The class writes the abstract body verbatim to a file and reads it back. The optional argument (the class's
% own heading) is given explicitly, so the class does not look ahead for one.
\begin{abstract}[\abstractname]
A cut-path between two vertices $u$ and $v$ of a graph is a set of edges that contains both a $u$--$v$ path and a cut separating $u$ from $v$.
It models two servers that communicate along the path while every other route between them uses one of the chosen edges.
We study the \MinCutPath{} problem, which asks for a cut-path with the fewest edges.
We prove that its decision version is NP-complete.
In graphs of diameter two and in graphs in which every two vertices can be separated by at most two edges, a minimum cut-path has $c(u,v) + d(u,v) - 1$ edges.
Here $c(u,v)$ is the size of a minimum cut and $d(u,v)$ is the distance.
In Erdős–Rényi random graphs above the connectivity threshold, the union of a minimum cut and a shortest path is within a factor $1+\epsilon$ of the optimum, for every fixed $\epsilon > 0$, with probability tending to one.

\end{abstract}

\begin{keywords}
NP-completeness \sep edge connectivity \sep graph diameter \sep Erdős–Rényi graphs \sep average-case approximation
\end{keywords}

\maketitle

% PDF author field: \maketitle of the class fills it with the author list followed by ", ", which the
% PDF shows as a stray character; set after \maketitle, before the first page is written.
\hypersetup{pdfauthor={Norbert Micheľ}}

\section{Introduction}
\label{sec:introduction}

Suppose that two trusted servers in a network must communicate while an adversary is denied every other route between them.
We model the network as a graph whose vertices are the servers and whose edges are the direct communication links.
The task is to select a set of edges that contains a connecting path, which carries the communication, and a separating cut, which guarantees that every other route between the two servers uses at least one of the selected edges.
We call such a set of edges a \emph{cut-path}.

Connectivity and separation are classical topics: multi-terminal minimum cuts~\cite{GomoryHu1961}, representations of all minimum cuts~\cite{NagamochiKameda1996}, and small cuts and edge connectivity~\cite{Mehlhorn2017Certifying} have been studied in detail, and paths and cuts are related by the max-flow min-cut duality~\cite{Cormen2022}.
These papers treat a path and a cut as separate objects.
A cut-path requires both in the same set of edges.


The \MinCutPath{} problem asks for a cut-path with the fewest edges between two given vertices.
To the best of our knowledge, the problem has not been studied before.
We settle its complexity and prove its basic structural properties.

Our contributions are the following.
First, we prove that the decision version of \MinCutPath{} is NP-complete, by a reduction from \ThreeSAT{} through an intermediate problem, \SSP.
Second, we prove that $\cp(u,v) = c(u,v) + d(u,v) - 1$ in graphs of diameter two and in graphs with minimum cut-value at most two.
Here $\cp(u,v)$ is the size of a minimum cut-path, $c(u,v)$ the size of a minimum cut, and $d(u,v)$ the distance; in both classes a minimum cut-path can be computed in polynomial time.
Third, dense Erdős–Rényi random graphs have diameter two with probability tending to one, so the formula applies to them.
For sparser random graphs, we show that the union of a minimum cut and a shortest path is within a factor of $1+\epsilon$ of the optimum on almost all inputs.

Section~\ref{sec:fundamentals} defines cut-paths and proves $\max\{c(u,v), d(u,v)\} \leq \cp(u,v) \leq c(u,v) + d(u,v) - 1$.
Section~\ref{sec:np-completeness} contains the NP-completeness proof.
Sections~\ref{sec:diameter-two} and~\ref{sec:cut-two} treat graphs of diameter two and graphs with minimum cut-value at most two, and Section~\ref{sec:random-graphs} treats random graphs.
Section~\ref{sec:conclusion} lists open problems.

\section{Fundamentals}
\label{sec:fundamentals}

Throughout the paper, we use standard notation from graph theory and computational complexity, following the conventions of~\cite{Diestel2025,Cormen2022}.
All graphs under consideration are undirected and finite, unless otherwise stated.
For a graph $G = (V, E)$, we let $n = |V|$ denote the number of vertices and $m = |E|$ the number of edges.

We often deal with two distinct vertices $u, v \in V$.
The value $d(u, v)$ denotes the distance between $u$ and $v$, i.e., the length of a shortest $u$--$v$ path in $G$.
Similarly, the value $c(u, v)$ denotes the minimum size of a cut that separates $u$ and $v$.
All problems studied in this paper are defined with respect to these two distinguished vertices $u$ and $v$, which are assumed to lie in the same connected component of $G$.

We identify a path with its set of edges.
For any path $P \subseteq E$, we denote by $G \setminus P$ the graph obtained from $G$ by removing all edges of $P$.
For a vertex $x \in V$ and a vertex subset $S \subseteq V$, we write $\deg(x, S)$ for the number of neighbors of $x$ that belong to $S$, that is, $\deg(x, S) = |\{\,y \in S : \{x, y\} \in E\,\}|$.



\begin{definition}[Cut-Path]
\label{def:cut-path}
Let $G = (V, E)$ be a graph, and let $u, v \in V$ be two distinct vertices.
A \emph{cut-path} between $u$ and $v$ is a set of edges $S \subseteq E$ for which there exist subsets $P, C \subseteq S$ such that:
    \begin{enumerate}
        \item $C$ is a cut separating $u$ and $v$ (i.e., removing $C$ disconnects $u$ from $v$),
        \item $P$ is a path connecting $u$ and $v$.
    \end{enumerate}
\end{definition}

Figure~\ref{fig:cut-path} shows a cut-path schematically.

\begin{figure}[H]
    \centering
    \includegraphics{cut-path.pdf}
    \caption{A cut-path between $u$ and $v$ (schematic)}
    \label{fig:cut-path}
\end{figure}

\begin{definition}[Cut-Paths]
\label{def:cut-paths}
Let $G = (V, E)$ be a graph, and let $u, v \in V$ be two distinct vertices.
The set of all cut-paths between $u$ and $v$, denoted by $\CP(u, v)$, is defined as:
\[
\CP(u, v) = \{ S \mid S \text{ is a cut-path between } u \text{ and } v \}.
\]
\end{definition}

\begin{definition}[Value of Minimum Cut-Path]
\label{def:cp-value}
Let $G = (V, E)$ be a graph, and let $u, v \in V$ be two distinct vertices.
The minimum size of a cut-path between $u$ and $v$, denoted by $\cp(u, v)$, is defined as:
\[
\cp(u, v) = \min\{ |S| \mid S \in \CP(u, v) \}.
\]
\end{definition}

A cut-path $S \in \CP(u, v)$ with $|S| = \cp(u, v)$ is called a \emph{minimum cut-path} between $u$ and $v$.
The corresponding optimization problem is stated as follows.

\begin{problem}[\MinCutPath{} (Optimization Version)]
Let $G=(V,E)$ be an undirected graph and $u,v \in V$ distinct vertices; cut-paths are defined in Definition~\ref{def:cut-path}.
The optimization version of the \MinCutPath{} problem is defined as follows:

\begin{center}
\begin{tabular}{p{0.18\linewidth}p{0.75\linewidth}}
\textbf{Input:} & A graph $G=(V,E)$ and vertices $u,v \in V$. \\
\textbf{Output:} & A minimum cut-path $S^* \subseteq E$ between $u$ and $v$, together with its cardinality
\[
\cp(u,v) = |S^*|.
\]
\end{tabular}
\end{center}
\end{problem}

The following elementary bounds relate $\cp(u, v)$ to the classical quantities $c(u, v)$ and $d(u, v)$; they are used repeatedly in the subsequent sections.

\begin{lemma}[Basic Bounds]
\label{lem:basic-bounds}
Let $G = (V, E)$ be a graph, and let $u, v \in V$ be two distinct vertices. Then
\[
\max\{c(u, v), d(u, v)\} \leq \cp(u, v) \leq c(u, v) + d(u, v) - 1.
\]
In particular, if $c(u, v) = 1$ or $d(u, v) = 1$, then $\cp(u, v) = c(u, v) + d(u, v) - 1$.
\end{lemma}

\begin{proof}
Every cut-path contains a $u$--$v$ cut, which has at least $c(u, v)$ edges, and a $u$--$v$ path, which has at least $d(u, v)$ edges; this gives the lower bound.
For the upper bound, let $C$ be a minimum $u$--$v$ cut and let $P$ be a shortest $u$--$v$ path.
The set $C \cup P$ is a cut-path.
Since removing $C$ disconnects $u$ from $v$, the path $P$ contains at least one edge of $C$, and hence $|C \cup P| \leq c(u, v) + d(u, v) - 1$.
If $c(u, v) = 1$ or $d(u, v) = 1$, the lower and the upper bound coincide.
\end{proof}


For random graphs, we use the following notion of approximation: a polynomial-time algorithm that is within a factor of $1+\epsilon$ of the optimum on almost all inputs.
Let $I(n)$ denote the set of all input instances of size $n$, and let $\OPT(x)$ denote the optimal value of the objective function for an instance $x \in I(n)$, i.e., the minimum for a minimization problem and the maximum for a maximization problem.

\begin{definition}[Average $(1+\epsilon)$-Approximation Scheme]
\label{def:approximation-scheme}
Let $\epsilon > 0$. An algorithm $A$ is called an \emph{average $(1+\epsilon)$-approximation scheme} if it meets the following criteria:
\begin{enumerate}
    \item $A$ is a polynomial-time algorithm, i.e., its running time is $O(P(n))$ for some polynomial $P(n)$.
    \item For each input size $n$, there exists a subset of inputs $I_{A}^{\mathrm{opt}}(n) \subseteq I(n)$ such that for every input $x \in I_{A}^{\mathrm{opt}}(n)$, the value $A(x)$ of the solution computed by $A$ satisfies
    \[
    \frac{A(x)}{\OPT(x)} < 1+\epsilon \quad \text{(for minimization problems)}
    \]
    \[
    \frac{\OPT(x)}{A(x)} < 1+\epsilon \quad \text{(for maximization problems)},
    \]
    where $\OPT(x)$ denotes the optimal value for the input $x$.
    \item In addition, the algorithm is average-case in the sense that
    \[
    \lim_{n \to \infty} \frac{\lvert I_{A}^{\mathrm{opt}}(n) \rvert}{\lvert I(n) \rvert} = 1.
    \]
\end{enumerate}
\end{definition}

\section{NP-completeness of the Min Cut-Path Problem}
\label{sec:np-completeness}
We prove that the decision version of the \MinCutPath{} problem is NP-complete, starting from the \ThreeSAT{} problem.





\begin{problem}[\ThreeSAT]
Let $U = \{x_1, x_2, \dots, x_n\}$ be a finite set of Boolean variables, and let
$\mathcal{C} = \{C_1, C_2, \dots, C_m\}$ be a collection of clauses, where each clause
$C_k$ is a disjunction of exactly three literals, and each literal is either a variable $x_i$ or its negation $\neg x_i$.

The \ThreeSAT{} decision problem is defined as follows:

\begin{center}
\begin{tabular}{p{0.12\linewidth}p{0.8\linewidth}}
\textbf{Input:} & A finite set of Boolean variables $U$ and a collection $\mathcal{C}$ of clauses, each containing exactly three literals.\\
\textbf{Question:} & Does there exist a truth assignment $\tau: U \rightarrow \{\mathsf{true}, \mathsf{false}\}$ such that every clause $C_k \in \mathcal{C}$ evaluates to \textsf{true} under $\tau$?
\end{tabular}
\end{center}
\end{problem}


\ThreeSAT{} is NP-complete~\cite{Cook1971Complexity}.
Throughout this section, $n$ and $m$ denote the numbers of variables and clauses of the \ThreeSAT{} instance under consideration.

The proof has two steps and uses an intermediate problem, \SSP.
First, we show that this problem is NP-complete by a polynomial-time reduction from \ThreeSAT{} (Section~\ref{sec:ssp}).
Second, we reduce \SSP{} to \MinCutPath{} (Section~\ref{sec:ssp-to-mcp}).


\subsection{The Separating Shortest Path Problem}\label{sec:ssp}
The \SSP{} problem asks whether the removal of some shortest $u$--$v$ path disconnects $u$ from $v$.

\begin{problem}[\SSP]
Let $G = (V,E)$ be an undirected graph and let $u,v \in V$ be two distinct vertices.
A shortest $u$--$v$ path $P$ is called a \emph{separating shortest path} if the graph $G \setminus P$ contains no $u$--$v$ path.
The \SSP{} decision problem is defined as follows:

\begin{center}
\begin{tabular}{p{0.12\linewidth}p{0.8\linewidth}}
\textbf{Input:} & A graph $G = (V,E)$ and two distinct vertices $u,v \in V$.\\
\textbf{Question:} & Does there exist a shortest $u$--$v$ path in $G$ that is a separating shortest path?
\end{tabular}
\end{center}
\end{problem}

Unlike \MinCutPath, \SSP{} is not an optimization problem, which makes the reduction from \ThreeSAT{} more direct.




\subsubsection{Gadgets and Basic Operations for the Reduction}

The reduction from \ThreeSAT{} to \SSP{} is built from \emph{gadgets} and \emph{basic operations} on them.

The construction consists of two parts:
\begin{enumerate}
    \item the \emph{base structure}, also referred to as the \emph{chain};
    \item the \emph{threads}, which are threaded through the base structure.
\end{enumerate}

Figure~\ref{fig:chain-threads} shows both parts.


\begin{figure}[H]
    \centering
    \includegraphics{chain-threads.pdf}
    \caption{Layout of the construction: the chain (black) and a thread (red)}
    \label{fig:chain-threads}
\end{figure}


\paragraph{Chain.}
The smallest gadget is the \emph{chain link}, a cycle with two distinguished vertices $p$ and $q$ that split it into two paths of equal length.


\begin{definition}[Chain Link]
\label{def:chain-link}
Let $G=(V,E)$ be a graph and let $p,q \in V$ be two distinct vertices.
A \emph{$(p,q)$-chain link} is a subgraph $L \subseteq G$ such that:
\begin{enumerate}
    \item $L$ is a simple cycle in $G$;
    \item the vertices $p$ and $q$ lie on $L$ and partition it into two internally disjoint $p$--$q$ paths;
    \item these two $p$--$q$ paths have equal length.
\end{enumerate}
We refer to one of the two $p$--$q$ paths as the \emph{positive path} of the chain link and to the other as the \emph{negative path}.
When a chain link is used inside a chain, its boundary vertices are exactly $p$ and~$q$, i.e., the points where the adjacent single edges of the chain are attached.
\end{definition}


\begin{figure}[H]
    \centering
    \includegraphics{chain-link.pdf}
    \caption{A chain link}
    \label{fig:chain-link}
\end{figure}



A chain alternates single edges and chain links; it begins and ends with a single edge.

\begin{definition}[Chain]
\label{def:chain}
Let $G=(V,E)$ be a graph, and let $u,v\in V$.
We say that $G$ is a \emph{$u$--$v$ chain} with $r$ chain links if it is the union of edge-disjoint subgraphs
\[
G = H_1 \cup L_1 \cup H_2 \cup L_2 \cup \cdots \cup L_r \cup H_{r+1},
\]
where:
\begin{enumerate}
    \item each $H_i$ consists of a single edge with end-vertices $p_i,q_i$;
    \item each $L_i$ is a $(q_i,p_{i+1})$-chain link;
    \item the first end-vertex $p_1$ coincides with $u$ and the last end-vertex $q_{r+1}$ coincides with $v$;
    \item two of these subgraphs share a vertex only if they are consecutive, and then they share exactly one vertex ($q_i$ or $p_{i+1}$).
\end{enumerate}
Every $u$--$v$ path in a chain traverses all single edges and, in each chain link, either its positive or its negative path.
\end{definition}

\begin{figure}[H]
    \centering
    \includegraphics{chain.pdf}
    \caption{A chain}
    \label{fig:chain}
\end{figure}



Let $C_1, C_2, \dots, C_m$ denote the clauses of the input formula and let $x_1, x_2, \dots, x_n$ denote its variables.
In the reduction, every chain link has one of three types:

\begin{enumerate}
    \item \textbf{Initialization chain link.}
    For each variable $x_i$ we introduce an initialization chain link, denoted $I_i$.
    This link marks the beginning of the part of the chain that corresponds to the variable $x_i$. The total number of initialization chain links is $n$.

    \item \textbf{Terminal chain link.}
    For each variable $x_i$ we also introduce a terminal chain link, denoted $T_i$.
    This link marks the end of the part of the chain associated with the variable $x_i$. The total number of terminal chain links is $n$.

    \item \textbf{Literal chain link.}
    For each clause $C_k = (\ell_1 \vee \ell_2 \vee \ell_3)$ and each $j \in \{1,2,3\}$, let $x_i$ be the variable of the literal $\ell_j$.
    We introduce a literal chain link, denoted $L_{j,k,i}$ (or $L_{j,k}$ for short).
    This link represents the occurrence of the variable $x_i$ as the $j$-th literal of the clause $C_k$.
    The total number of literal chain links is $3m$.

\end{enumerate}

The chain links that belong to one variable are consecutive in the chain, as illustrated in Figure~\ref{fig:chain-link-types}.

\begin{figure}[H]
    \centering
    \includegraphics{chain-link-types.pdf}
    \caption{Chain links of a variable $x_i$: the initialization link, the literal links, and the terminal link}
    \label{fig:chain-link-types}
\end{figure}

In total, the construction contains $3m + 2n$ chain links.

\paragraph{Threads and threading.}
\emph{Threads} are paths that pass through the chain links; inserting such a path into the chain is called \emph{threading}.

\begin{definition}[Thread]
\label{def:thread}
Let $G=(V,E)$ be a graph containing a $u$--$v$ chain with chain links $L_1,L_2,\dots,L_r$.
A \emph{thread} is a $u$--$v$ path $P$ in $G$ such that:
\begin{enumerate}
    \item for each chain link $L_i$, the path $P$ shares at most one edge with $L_i$, and $P$ contains none of the single edges of the chain;
    \item $P$ is edge-disjoint from every other thread and shares only the vertices $u$ and $v$ with it.
\end{enumerate}
\end{definition}


Threading subdivides an edge of a chain link that no thread uses yet into three edges and routes the thread through the middle one.

\begin{definition}[Threading]
\label{def:threading}
Let $L$ be a chain link in a $u$--$v$ chain, and let $P$ be a thread not yet passing through $L$.
The \emph{threading operation} on $(P,L)$ consists of the following elementary steps:
\begin{enumerate}
    \item Select an edge $e=\{x,y\}$ of $L$ that is not used by any existing thread.
    \item Subdivide $e$ into a path $x$--$z_1$--$z_2$--$y$ of three edges, introducing two new vertices $z_1,z_2$.
    \item Extend $P$ so that it traverses the middle edge $\{z_1,z_2\}$, called the \emph{crossing edge}, as its unique crossing of the chain link $L$.
\end{enumerate}
If the selected edge $e$ belongs to the positive path of $L$, we say that $P$ is \emph{positively threaded through $L$}; if $e$ belongs to the negative path of $L$, we say that $P$ is \emph{negatively threaded through $L$}.
In general, we say that the thread $P$ is \emph{threaded through $L$ at the edge $e$}.
\end{definition}

\begin{figure}[H]
    \centering
    \includegraphics{threading.pdf}
    \caption{The threading operation}
    \label{fig:threading}
\end{figure}


Every thread has one of three types:

\begin{enumerate}
    \item \textbf{2-synchronization thread.}
    This thread passes through an initialization chain link $I_i$ and the corresponding terminal chain link $T_i$ (Figure~\ref{fig:two-synchronization}).
    Threads of this type force every separating shortest path to traverse paths of the same sign in $I_i$ and $T_i$.

    \begin{figure}[H]
        \centering
        \includegraphics{two-synchronization-threads.pdf}
        \caption{The two 2-synchronization threads of a variable $x_i$}
        \label{fig:two-synchronization}
    \end{figure}

    \item \textbf{3-synchronization thread.}
    This thread passes through the initialization chain link $I_i$, a literal chain link $L_{j,k,i}$, and the terminal chain link $T_i$ corresponding to the same variable $x_i$ (Figure~\ref{fig:three-synchronization}).
    Threads of this type force every separating shortest path to traverse the path of this sign in $L_{j,k,i}$ as well.

    \begin{figure}[H]
        \centering
        \includegraphics{three-synchronization-threads.pdf}
        \caption{The two 3-synchronization threads of a literal chain link $L_{j,k,i}$}
        \label{fig:three-synchronization}
    \end{figure}

    \item \textbf{Clause thread.}
    For each clause $C_k = (\ell_{1} \vee \ell_{2} \vee \ell_{3})$, a clause thread is introduced that passes through the three literal chain links $L_{1,k}$, $L_{2,k}$, $L_{3,k}$ corresponding to $\ell_{1}, \ell_{2}, \ell_{3}$ (Figure~\ref{fig:clause-thread}).
    It encodes the clause: a separating shortest path exists only if at least one of the three literals is satisfied.

    \begin{figure}[H]
        \centering
        \includegraphics{clause-thread.pdf}
        \caption{The clause thread of a clause $C_k = (x_1 \vee \neg x_2 \vee x_3)$}
        \label{fig:clause-thread}
    \end{figure}
\end{enumerate}

\subsubsection{NP-completeness of \textsc{Separating Shortest Path}}
We prove that \SSP{} is NP-complete by a polynomial-time reduction from \ThreeSAT{} that uses the chain and the threads.

\paragraph{Auxiliary procedures.}
The reduction uses the following auxiliary procedures:
\begin{itemize}
    \item $\textsc{BuildChain}(u,v,r)$: constructs a $u$--$v$ chain with $r$ chain links, where each chain link is realized as a $4$-cycle.
    \item $\textsc{Thread}(u,v,[(K_1,\sigma_1),\dots,(K_t,\sigma_t)])$: creates a new thread from $u$ to $v$ that passes through the chain links $K_1,\dots,K_t$ in this order, where $\sigma_i \in \{+,-\}$ specifies whether the thread is threaded through $K_i$ positively or negatively.
    The $t$ crossing edges are joined to each other and to the vertices $u$ and $v$ by $t+1$ new, internally vertex-disjoint paths, which we call the \emph{connecting paths} of the thread.
    \item $\textsc{Calibrate}(G)$: equalizes the lengths of the relevant paths in two steps.
    First, in every chain link whose positive and negative paths have different lengths, an edge of the shorter path that is not a crossing edge is subdivided repeatedly until both paths have the same length.
    Second, let $\Lambda$ denote the number of edges of a $u$--$v$ path in the resulting chain; every connecting path of every thread is subdivided until it has at least $\Lambda$ edges.
\end{itemize}

\paragraph{Auxiliary notation.}
We also introduce the following notational conventions:
\begin{itemize}
    \item Occurrences: for each variable $x_i$, we denote by $O_i$ the set of all pairs $(j,k)$ such that the $j$-th literal of the clause $C_k$ is $x_i$ or $\neg x_i$.
    \item Sign function: for each literal $\ell$ we define
\[
s(\ell) =
\begin{cases}
+ & \text{if $\ell = x_i$ for some variable $x_i$,}\\
- & \text{if $\ell = \neg x_i$ for some variable $x_i$.}
\end{cases}
\]
\end{itemize}

\begin{theorem}
\label{thm:ssp-np-complete}
The \SSP{} problem is NP-complete.
\end{theorem}

\begin{proof}
We first describe the reduction, then prove its correctness, and finally bound its running time and show that the problem belongs to NP.

\paragraph{Reduction from 3-SAT.}
Given an instance of \ThreeSAT, we describe how to construct a graph $G$ with two distinguished vertices $u,v$ such that $G$ contains a separating shortest $u$--$v$ path if and only if the given formula is satisfiable.
The construction is summarized in Algorithm~\ref{alg:reduction}.


\begin{algorithm}[H]
\caption{Reduction from \ThreeSAT{} to \SSP}
\label{alg:reduction}
\begin{algorithmic}[1]
\Require Variables $U=\{x_1,\dots,x_n\}$, clauses $\mathcal{C}=\{C_1,\dots,C_m\}$
\Ensure Graph $G$, vertices $u,v$
\State $G \gets \textsc{BuildChain}(u,v,2n+3m)$
\For{$i \gets 1$ to $n$}
    \State reserve the next $2+|O_i|$ free chain links
    \State label the first of them as $I_i$ and the last of them as $T_i$
    \For{each $(j,k) \in O_i$}
        \State label the next free reserved link as $L_{j,k,i}$
    \EndFor
\EndFor
\For{$i \gets 1$ to $n$}
    \State $\textsc{Thread}(u,v,[(I_i,+),(T_i,-)])$
    \State $\textsc{Thread}(u,v,[(I_i,-),(T_i,+)])$
    \For{each $(j,k) \in O_i$}
        \State $\textsc{Thread}(u,v,[(I_i,+),(L_{j,k,i},-),(T_i,+)])$
        \State $\textsc{Thread}(u,v,[(I_i,-),(L_{j,k,i},+),(T_i,-)])$
    \EndFor
\EndFor
\For{$k \gets 1$ to $m$}
    \State $\ell_1,\ell_2,\ell_3 \gets$ the literals of $C_k$
    \State $\textsc{Thread}(u,v,[(L_{1,k},s(\ell_1)),(L_{2,k},s(\ell_2)),(L_{3,k},s(\ell_3))])$
\EndFor
\State $\textsc{Calibrate}(G)$
\State \Return $G$, $u$, $v$
\end{algorithmic}
\end{algorithm}

In words, the chain consists of $n$ consecutive blocks, one for each variable; the block of $x_i$ starts with $I_i$, ends with $T_i$, and contains the literal chain links of all occurrences of $x_i$ in between.
Each variable receives a pair of 2-synchronization threads, each occurrence of a variable receives a pair of 3-synchronization threads, and each clause receives one clause thread.
The final calibration step guarantees that both paths of every chain link have the same length and that the connecting paths of the threads are sufficiently long.

\paragraph{Correctness.}
We now show that the constructed instance of \SSP{} is equivalent to the original \ThreeSAT{} instance.
We call the $u$--$v$ paths of the chain \emph{chain paths}; a chain path is determined by choosing, in every chain link, either its positive or its negative path.
We say that a chain path $P$ \emph{hits} a thread if $P$ and the thread share an edge.
Since a thread shares only its crossing edges with the chain, a chain path $P$ hits a thread created by $\textsc{Thread}(u,v,[(K_1,\sigma_1),\dots,(K_t,\sigma_t)])$ if and only if, for some $i$, the path $P$ traverses the path of $K_i$ with the sign $\sigma_i$.

\paragraph{Shortest paths.}
After the calibration, all chain paths have the same length $\Lambda$.
Every other $u$--$v$ path $Q$ in $G$ contains an edge of some connecting path.
Since the internal vertices of a connecting path have degree two, $Q$ contains the whole connecting path, which has at least $\Lambda$ edges, and at least one further edge, because no connecting path joins $u$ and $v$ directly.
Hence $|Q| > \Lambda$, and the shortest $u$--$v$ paths in $G$ are exactly the chain paths.

\paragraph{Synchronization of variable gadgets.}
Every thread is a $u$--$v$ path. Therefore, a chain path $P$ can be separating only if it hits every thread.
\begin{itemize}
    \item \emph{2-synchronization threads:} the chain path $P$ hits both 2-synchronization threads of the variable $x_i$ (Figure~\ref{fig:two-synchronization}) if and only if it traverses paths of the same sign in $I_i$ and in $T_i$.
    Indeed, the first thread requires the positive path in $I_i$ or the negative path in $T_i$, and the second thread requires the negative path in $I_i$ or the positive path in $T_i$.
    \item \emph{3-synchronization threads:} if $P$ traverses paths of the same sign in $I_i$ and $T_i$, then it hits both 3-synchronization threads of a literal link $L_{j,k,i}$ (Figure~\ref{fig:three-synchronization}) if and only if it traverses the path of this sign in $L_{j,k,i}$ as well.
    Indeed, one of the two threads is hit exactly when $P$ uses the positive path in $I_i$ or $T_i$ or the negative path in $L_{j,k,i}$, and the other one exactly when $P$ uses the negative path in $I_i$ or $T_i$ or the positive path in $L_{j,k,i}$.
\end{itemize}
Consequently, a chain path hits all synchronization threads if and only if, for every variable $x_i$, it traverses paths of the same sign in all chain links of the block of $x_i$.
We call such a chain path \emph{consistent}.
Consistent chain paths correspond to truth assignments: we set $\tau(x_i)=\mathsf{true}$ if the path traverses the positive paths in the block of $x_i$, and $\tau(x_i)=\mathsf{false}$ otherwise.

\paragraph{Clause threads.}
Let $P$ be a consistent chain path with the truth assignment $\tau$, and let $C_k = (\ell_1 \vee \ell_2 \vee \ell_3)$ be a clause.
The path $P$ hits the clause thread of $C_k$ if and only if, for some $j$, it traverses the path of $L_{j,k}$ with the sign $s(\ell_j)$, i.e., if and only if $\tau$ satisfies the literal $\ell_j$.
Thus a consistent chain path hits all clause threads if and only if its truth assignment satisfies every clause.

\paragraph{Two-way correspondence.}
We complete the proof by establishing the equivalence:
\begin{enumerate}[(a)]
    \item If there exists a separating shortest path $P$ in $G$, then one can extract a satisfying assignment for the original \ThreeSAT{} instance.

    Since $P$ is a shortest $u$--$v$ path, it is a chain path, and since it is separating, it hits every thread.
    By the synchronization property, $P$ is consistent, so it determines a truth assignment $\tau$.
    Since $P$ also hits every clause thread, $\tau$ satisfies at least one literal in every clause.
    Thus, $\tau$ is a satisfying assignment for the original formula.

    \item If there exists a satisfying assignment for the \ThreeSAT{} instance, then one can construct a separating shortest path in $G$.

Given a satisfying assignment $\tau$, let $P$ be the consistent chain path that corresponds to $\tau$: in all chain links of the block of $x_i$, it traverses the positive path if $\tau(x_i)=\mathsf{true}$ and the negative path if $\tau(x_i)=\mathsf{false}$.
The path $P$ is a shortest $u$--$v$ path, and it hits every thread.

It remains to show that $P$ is separating, i.e., that $G\setminus P$ contains no $u$--$v$ path.
In $G \setminus P$, all single edges of the chain are missing, and in every chain link $K$ only the path not traversed by $P$ remains; we call it the \emph{unused path} of $K$.
A crossing edge that lies on $P$ is removed together with the two edges of the chain link adjacent to it; hence the two connecting paths that end at this crossing edge are dead ends in $G \setminus P$.
A crossing edge on an unused path remains and joins its two adjacent connecting paths to this unused path.
Consequently, a path in $G \setminus P$ can move between connecting paths and unused paths only at crossing edges that are not on $P$.
We analyze all ways in which a path starting at $u$ could progress in $G \setminus P$:
\begin{enumerate}
  \item \emph{Using a chain edge.} The single edge of the chain incident to $u$ lies on $P$, so this option is blocked immediately.
  \item \emph{Using a synchronization thread.} Such a thread first reaches $I_i$.
  If it crosses the path of $I_i$ traversed by $P$, it ends there.
  Otherwise, it reaches the unused path of $I_i$.
  The threads that cross this unused path are one 2-synchronization thread, whose next crossing edge is in $T_i$, and one 3-synchronization thread for each literal link $L_{j,k,i}$, whose next crossing edge is in $L_{j,k,i}$; all these next crossing edges lie on $P$.
  Hence every continuation from the unused path of $I_i$ leads back to $u$ or to a dead end.
  \item \emph{Using a clause thread.} The clause thread of $C_k$ passes through the literal links $L_{1,k}$, $L_{2,k}$, $L_{3,k}$.
  Since $\tau$ satisfies $C_k$, at least one of its crossing edges lies on $P$, so the clause thread itself is interrupted before it reaches $v$.
  At a crossing edge that is not on $P$, a path can leave the clause thread and enter the unused path of the literal link $L_{j,k,i}$.
  Apart from the clause thread, this unused path is crossed only by a 3-synchronization thread, and the two connecting paths of this thread that start there lead to crossing edges in $I_i$ and $T_i$ that lie on $P$, i.e., to dead ends.
\end{enumerate}
Since all possible continuations from $u$ are blocked in $G\setminus P$, there is no $u$--$v$ path after removing $P$; hence $P$ is separating.
\end{enumerate}

\paragraph{Running time of the reduction.}
It remains to show that the reduction runs in polynomial time.
The chain has $2n+3m$ chain links, and the construction creates $2n+7m$ threads, each of which has at most three crossing edges and at most four connecting paths.
Before the calibration, the chain therefore has $O(n+m)$ edges.
Balancing the two paths of a chain link adds at most two edges for each of its crossing edges, so $\Lambda = O(n+m)$.
Finally, each of the $O(n+m)$ connecting paths is subdivided until it has $\Lambda$ edges, which requires $O((n+m)^2)$ edges in total.
Hence, the size of $G$ and the running time of Algorithm~\ref{alg:reduction} are polynomial in the input size.

\paragraph{Membership in NP.}
A candidate path $P$ is verified in polynomial time: check that $P$ is a shortest $u$--$v$ path, remove its edges from $G$, and test by breadth-first search whether a $u$--$v$ path remains.
Hence \SSP{} belongs to NP.
\end{proof}

\subsection{From Separating Shortest Path to Min Cut-Path}\label{sec:ssp-to-mcp}
We now reduce \SSP{} to \MinCutPath, which we state in decision form.

\begin{problem}[\MinCutPath{} (Decision Version)]
Let $G=(V,E)$ be an undirected graph and $u,v \in V$ distinct vertices; cut-paths are defined in Definition~\ref{def:cut-path}.
The decision version of the \MinCutPath{} problem is defined as follows:

\begin{center}
\begin{tabular}{p{0.12\linewidth}p{0.8\linewidth}}
\textbf{Input:} & A graph $G=(V,E)$, vertices $u,v \in V$, and an integer $k$. \\
\textbf{Question:} & Does there exist a cut-path $F \subseteq E$ between $u$ and $v$ with $|F|\leq k$?
\end{tabular}
\end{center}
\end{problem}

The decision version is NP-complete, and hence the optimization version is NP-hard.

\begin{theorem}
\label{thm:mcp-np-complete}
The decision version of the \MinCutPath{} problem is NP-complete.
\end{theorem}

\begin{proof}
\textbf{Reduction.}
Let $G$ be an instance of \SSP{} with distinguished vertices $u,v$.
We may assume that $u$ and $v$ lie in the same connected component of $G$, since otherwise the answer is negative.
We define $G' := G$ and set the threshold $k := d_G(u,v)$, the length of a shortest $u$--$v$ path in $G$.
We then ask whether $G'$ admits a cut-path of size at most $k$.

\textbf{Correctness.}
Suppose that $G$ admits a separating shortest path $P$.
Then $|P| = d_G(u,v) = k$, and $P$ itself is a cut-path of size $k$, since (i) it is a $u$--$v$ path and (ii) its removal disconnects $u$ and $v$, so $P$ is also a $u$--$v$ cut.
Conversely, suppose that $G'$ has a cut-path $F$ with $|F| \leq k$.
Then $F$ contains a $u$--$v$ path $P$, which has at least $d_G(u,v) = k$ edges; hence $F = P$ and $P$ is a shortest $u$--$v$ path in $G$.
Since $F$ also contains a $u$--$v$ cut, $G \setminus P$ contains no $u$--$v$ path.
Hence $P$ is a separating shortest path, and the two instances are equivalent.

\textbf{Membership in NP.}
A candidate cut-path $F$ can be verified in polynomial time:
check that $|F| \leq k$, that $F$ contains a $u$--$v$ path, and that $G\setminus F$ contains no $u$--$v$ path (e.g., by BFS or DFS).
All conditions can be tested in polynomial time.

\textbf{Conclusion.}
By Theorem~\ref{thm:ssp-np-complete}, \SSP{} is NP-complete, and it reduces to \MinCutPath{} in polynomial time.
Since \MinCutPath{} is in NP (by polynomial verification), it follows that \MinCutPath{} is NP-complete.
\end{proof}

\section{Graphs of Diameter Two}
\label{sec:diameter-two}



The first class in which the problem can be solved efficiently consists of the graphs of diameter two, a class studied in algebraic graph theory~\cite{GodsilRoyle2001}.

\begin{definition}[Diameter]
\label{def:diameter}
A graph $G = (V, E)$ has \emph{diameter at most $k$} if $d(x, y) \leq k$ for every pair of vertices $x, y \in V$.
The smallest such $k$ is the \emph{diameter} of $G$, denoted by $\diam(G)$.
\end{definition}

We derive a formula for $\cp(u,v)$ in graphs of diameter two; it yields a polynomial-time algorithm (Remark~\ref{rem:algorithm}).
The proof uses a decomposition: in any graph, a cut $C$ induces a partition of the vertex set into four subsets $I, J, K, L$.



Each vertex belongs to one of the two sides $A_1, A_2$ of the cut, and it either is incident to an edge of the cut or not:

\begin{lemma}
\label{lem:cut-decomposition}
Let $G = (V, E)$ be a graph, and let $C$ be a cut that divides $G$ into two sets $A_1$ and $A_2$ (i.e., $V = A_1 \cup A_2$, $A_1 \cap A_2 = \emptyset$, and $C$ is the set of all edges between $A_1$ and $A_2$).
Then, there exists a unique decomposition of $A_1$ and $A_2$ into subsets $I, J, K, L$ such that:
\begin{itemize}
    \item $I \cup J = A_1$ and $I \cap J = \emptyset$,
    \item $K \cup L = A_2$ and $K \cap L = \emptyset$,
    \item $\forall x \in J \cup K, \, \exists y \in V: \{x, y\} \in C,$
    \item $\forall x \in I \cup L, \, \nexists y \in V: \{x, y\} \in C.$
\end{itemize}
\end{lemma}


In a graph of diameter two, $I$ or $L$ is empty: otherwise a vertex of $I$ would be at distance at least three from a vertex of $L$, since every path between them passes through $J$ and $K$.


\begin{lemma}
\label{lem:empty-i-or-l}
Let $G = (V, E)$ be a graph of diameter two, let $C$ be a cut that divides $G$ into two sets $A_1$ and $A_2$, and let $I, J, K, L$ be the decomposition of $A_1$ and $A_2$ described in Lemma~\ref{lem:cut-decomposition}.
Then at least one of the sets $I$ and $L$ is empty:
\[
I = \emptyset \quad\text{or}\quad L = \emptyset.
\]
\end{lemma}

We also need the following general fact about the intersection of a cut and a path.

\begin{lemma}
\label{lem:odd-intersection}
Let $G = (V, E)$ be a graph, and let $u, v \in V$ be two distinct vertices.
Let $C$ be a cut that divides $G$ into two sets $A_1$ and $A_2$ with $u \in A_1$ and $v \in A_2$, and let $P$ be a path that connects $u$ and $v$.
Then the path $P$ contains an odd number of edges of $C$:
\[
|C \cap P| \equiv 1 \pmod{2}.
\]
\end{lemma}

\begin{proof}
Let $S = C \cap P$ be the set of edges in the intersection of the cut and the path.
When we traverse the path $P$ from $u$ to $v$, each edge of $S$ changes the side of the cut (i.e., it switches between $A_1$ and $A_2$), while every other edge of $P$ stays on the same side.
The path starts in $A_1$ and ends in $A_2$.
If the path crossed the cut an even number of times, its first and its last vertex would lie on the same side, which contradicts the assumption that $u \in A_1$ and $v \in A_2$.
Therefore, the path crosses the cut an odd number of times, i.e., $|S|$ is odd.
\end{proof}

Lemmas~\ref{lem:empty-i-or-l} and~\ref{lem:odd-intersection} give the value of $\cp(u,v)$ in graphs of diameter two.

\begin{theorem}
\label{thm:diameter-two}
Let $G = (V, E)$ be a connected graph of diameter two, and let $u, v \in V$ be two distinct vertices.
Then the value of the minimum cut-path between $u$ and $v$ is
\[
\cp(u, v) = c(u, v) + d(u, v) - 1.
\]
\end{theorem}


\begin{proof}

We distinguish two cases based on the values of $d(u, v)$ and $c(u, v)$.

\begin{enumerate}
    \item \textbf{Case} $d(u, v) = 1$ \textbf{or} $c(u, v) = 1$:

    By Lemma~\ref{lem:basic-bounds}, in this case we have
    \[
    \cp(u, v) = c(u, v) + d(u, v) - 1.
    \]

    \item \textbf{Case} $d(u, v) = 2$ \textbf{and} $c(u, v) \geq 2$:

    By Lemma~\ref{lem:basic-bounds},
    \[
    c(u, v) \leq \cp(u, v) \leq c(u, v) + 2 - 1 = c(u, v) + 1.
    \]

Assume for contradiction that $\cp(u, v) \neq c(u, v)+1$.
Then $\cp(u, v) = c(u, v)$, i.e., a minimum cut-path $S$ has the same size as a minimum cut.
Since $S$ contains a $u$--$v$ cut, the set $S$ itself is a minimum cut, which we denote by $C$, and the $u$--$v$ path $P$ contained in $S$ is a subset of this cut.

By Lemma~\ref{lem:odd-intersection}, the number of edges of $P$ is odd, and since $d(u,v) = 2$, the path $P$ has length at least 3.
Let $I, J, K, L$ be the decomposition induced by the cut $C$.
By Lemma~\ref{lem:empty-i-or-l} and by the symmetry between $u$ and $v$, we may assume that $L = \emptyset$ and $u \in K$, so the graph is partitioned into the three sets $I, J, K$.
We observe that $v \in J$: the last edge of $P$ belongs to the cut and is incident to $v$, so the vertex $v$ cannot belong to the set $I$.
This structure is illustrated in Figure~\ref{fig:diameter-two-structure}.

\begin{figure}[H]
    \centering
    \includegraphics[width=60mm]{diameter-two-structure.png}
    \caption{Structure of the cut $C$ in the proof of Theorem~\ref{thm:diameter-two}; the edges counted in the estimate are drawn in red}
    \label{fig:diameter-two-structure}
\end{figure}



We now estimate the size of the cut $C$:
\[
\cp(u, v) = c(u,v) = |C| \geq \deg(u, I) + \deg(u, J) + 2 + (|K| - 2),
\]
where the estimate is obtained by counting the edges of the cut according to their end-vertex in $K$:
\begin{itemize}
    \item $\deg(u, I)$ and $\deg(u, J)$ denote the numbers of edges that lead from the vertex $u$ to the opposite side of the minimum cut,

    \item the path $P$ is contained in the cut and has length at least three, so its edges alternate between the two sides of the cut; hence its third vertex is a vertex of $K$ other than $u$ that contributes at least two edges to the cut,

    \item each of the remaining $|K| - 2$ vertices of the set $K$ contributes at least one edge to the cut.
\end{itemize}


The edges counted on the right-hand side are drawn in red in Figure~\ref{fig:diameter-two-structure}.

Since $\deg(u, I) \geq 0$ and $\deg(u, K) \leq |K| - 1$, we obtain
\[
\deg(u, I) + \deg(u, J) + 2 + (|K| - 2) \geq
0 + \deg(u, J) + \deg(u, K) + 1 = \deg(u) + 1 \geq c(u, v) + 1,
\]
where the equality holds because $u$ has no neighbors in $I$, and the last inequality holds because the edges incident to $u$ form a $u$--$v$ cut.
This contradicts the assumption that $\cp(u, v) = c(u, v)$; hence $\cp(u, v) = c(u, v) + 1 = c(u, v) + d(u, v) - 1$.\qedhere
\end{enumerate}
\end{proof}

\begin{remark}
\label{rem:algorithm}
Whenever $\cp(u, v) = c(u, v) + d(u, v) - 1$, the proof of Lemma~\ref{lem:basic-bounds} shows that the union of an arbitrary minimum $u$--$v$ cut and an arbitrary shortest $u$--$v$ path is a minimum cut-path.
Hence, in graphs of diameter two, a minimum cut-path can be found by one minimum-cut computation and one breadth-first search.
\end{remark}

\section{Graphs with Cut-Value at Most Two}
\label{sec:cut-two}

The same formula holds in graphs in which the minimum cut-value of every pair of vertices is at most two.

Such graphs have the structure of a \emph{cactus graph} (cf.~\cite{Mehlhorn2017Certifying}): they consist of cycles and single edges, any two of which share at most one vertex, an articulation point.
Indeed, if two distinct cycles shared an edge, their union would contain two vertices joined by three edge-disjoint paths; separating these two vertices would require at least three edges, contradicting the assumption $c(x,y) \leq 2$.

\begin{theorem}
\label{thm:cut-two}
Let $G = (V, E)$ be a connected graph such that $c(x, y) \leq 2$ for every pair of distinct vertices $x, y \in V$.
Then, for any two distinct vertices $u,v \in V$, the value of the minimum cut-path is
\[
\cp(u, v) = c(u, v) + d(u, v) - 1.
\]
\end{theorem}

\begin{proof}
We distinguish two cases based on the values of $d(u, v)$ and $c(u, v)$.

\begin{enumerate}
    \item \textbf{Case} $d(u, v) = 1$ \textbf{or} $c(u, v) = 1$:

    The formula holds by Lemma~\ref{lem:basic-bounds}.

    \item \textbf{Case} $d(u, v) \geq 2$ \textbf{and} $c(u, v) = 2$:

    By Lemma~\ref{lem:basic-bounds}, we have
    \[
    d(u, v) \leq \cp(u, v) \leq d(u, v) + 1.
    \]
    Assume for contradiction that $\cp(u, v) = d(u, v)$.
    Then a minimum cut-path is a shortest $u$--$v$ path $P$ that contains a $u$--$v$ cut, so the graph $G \setminus P$ contains no $u$--$v$ path.

    Since $c(u, v) = 2$, no single edge separates $u$ from $v$; hence no edge of $P$ is a bridge, and every edge of $P$ lies on a cycle of the cactus graph $G$.
    The path $P$ therefore passes through a sequence of cycles $C_1, C_2, \dots, C_t$, and in each cycle $C_i$ it follows one of the two arcs between the vertex where it enters $C_i$ and the vertex where it leaves $C_i$.
    The other arc of each cycle $C_i$ is edge-disjoint from $P$, and the concatenation of these arcs connects $u$ and $v$ in $G \setminus P$.
    This is a contradiction.
    Therefore,
    \[
    \cp(u, v) = d(u, v) + 1 = c(u, v) + d(u, v) - 1. \qedhere
    \]
\end{enumerate}
\end{proof}

By Remark~\ref{rem:algorithm}, a minimum cut-path in such a graph is again obtained as the union of a minimum cut and a shortest path.

\section{Erdős–Rényi Graphs}
\label{sec:random-graphs}

We study the \MinCutPath{} problem in Erdős–Rényi random graphs $G(n, p)$, in which each of the possible edges on $n$ vertices is present independently with probability $p$.
Throughout this section, $\log$ denotes the natural logarithm.

Dense random graphs, i.e., those with a constant edge probability $p$, have diameter two with probability tending to one~\cite{Bollobas2001}.

\begin{theorem}[\cite{Bollobas2001}]
\label{thm:random-diameter-two}
Let $G(n, p)$ be an Erdős–Rényi random graph with $n$ vertices and edge probability $p = \frac{1}{\alpha}$, where $\alpha > 1$ is a constant.
Then the probability that $G(n, p)$ has diameter two converges to 1 as $n \to \infty$:
\[
\lim_{n \to \infty} \mathbb{P}\!\left[\diam\bigl(G\bigl(n, \tfrac{1}{\alpha}\bigr)\bigr) = 2\right] = 1.
\]
\end{theorem}

By Theorem~\ref{thm:diameter-two} and Remark~\ref{rem:algorithm}, a minimum cut-path in a graph of diameter two has $c(u, v) + d(u, v) - 1$ edges and can be computed in polynomial time.
Consequently, for dense random graphs, the \MinCutPath{} problem can be solved exactly in polynomial time on almost all inputs.


Consider now sparse Erdős–Rényi random graphs $G(n, p)$ with $p = \frac{\alpha \log n}{n}$, where $\alpha > 1$ is a constant.
We state two known results on the diameter and the edge connectivity of such graphs~\cite{Bollobas2001,Frieze2016} and prove two consequences for the degrees and for the values $c(u,v)$.


\begin{theorem}[Diameter~\cite{Bollobas2001}]
\label{thm:random-diameter}
Let $G(n, p)$ be an Erdős–Rényi random graph with $p = \frac{\alpha \log n}{n}$, where $\alpha > 1$. Then:
\[
\lim_{n\to\infty} \mathbb{P}\left[ \diam(G) \leq \frac{\log n}{\log \log n} \right] = 1.
\]
\end{theorem}


\begin{theorem}[Connectivity~\cite{Frieze2016}]
\label{thm:random-connectivity}
Let $G(n, p)$ be an Erdős–Rényi random graph with $p = \frac{\alpha \log n}{n}$, where $\alpha > 1$. Then:
\[
\lim_{n\to\infty} \mathbb{P}\left[ G \text{ is } k\text{-edge-connected} \right] = 1,
\]
where $k = \delta(G)$ denotes the minimum degree of $G$.
\end{theorem}



\begin{lemma}[Degree Bounds]
\label{lem:degree-bounds}
Let $G(n, p)$ be an Erdős–Rényi random graph with $p = \frac{\alpha \log n}{n}$, where $\alpha > 1$ is a fixed constant.
Then there exist constants $0 < \beta_1 < \beta_2$, depending only on $\alpha$, such that
\[
\lim_{n\to\infty} \mathbb{P}\Bigl( \forall v \in V(G): \ \beta_1 \log n \leq \deg(v) \leq \beta_2 \log n \Bigr) = 1.
\]
\end{lemma}

\begin{proof}
The degree of a fixed vertex $v$ has the binomial distribution with mean $\mu = (n-1)p = (1-o(1))\,\alpha \log n$.
For $a > 0$, let $h(a) = a \log a - a + 1$.
By the Chernoff bounds (see, e.g.,~\cite{Frieze2016}),
\[
\mathbb{P}\bigl[\deg(v) \leq a\mu\bigr] \leq e^{-\mu h(a)} \ \text{ for } 0 < a < 1
\qquad\text{and}\qquad
\mathbb{P}\bigl[\deg(v) \geq a\mu\bigr] \leq e^{-\mu h(a)} \ \text{ for } a > 1.
\]
Since $h(a) \to 1$ as $a \to 0$, $h(a) \to \infty$ as $a \to \infty$, and $\alpha > 1$, we can fix constants $0 < a_1 < 1 < a_2$ such that $\alpha h(a_1) > 1$ and $\alpha h(a_2) > 1$.
For these constants, both probabilities above are $o(1/n)$, and the union bound over all $n$ vertices shows that, with probability tending to one, every degree lies between $a_1 \mu$ and $a_2 \mu$.
The claim follows, for instance, with $\beta_1 = a_1 \alpha / 2$ and $\beta_2 = 2 a_2 \alpha$.
\end{proof}


The next lemma bounds the values $c(u,v)$.

\begin{lemma}
\label{lem:connectivity-bounds}
Let $G(n, p)$ be an Erdős–Rényi random graph with $p = \frac{\alpha \log n}{n}$, where $\alpha > 1$, and let $\beta_1, \beta_2$ be the constants from Lemma~\ref{lem:degree-bounds}.
Then
\[
\lim_{n\to\infty} \mathbb{P}\Bigl( \forall u \neq v: \ \beta_1 \log n \leq c(u,v) \leq \beta_2 \log n \Bigr) = 1.
\]
\end{lemma}


\begin{proof}
By Theorem~\ref{thm:random-connectivity} and Lemma~\ref{lem:degree-bounds}, with probability tending to one, the graph $G(n, p)$ is $k$-edge-connected for $k = \delta(G)$ and all its degrees lie between $\beta_1 \log n$ and $\beta_2 \log n$.
In this case, all pairs of distinct vertices $u, v$ satisfy
\[
c(u, v) \geq \min_{x \neq y} c(x, y) \geq \min_{z \in V(G)} \deg(z) \geq \beta_1 \log n.
\]


On the other hand, the edges incident to $u$ form a cut that separates $u$ from $v$.
This cut need not be a minimum one, but it gives the upper bound
\[
c(u,v)\leq \deg(u) \leq \beta_2 \log n.
\]

Combining these two inequalities, we obtain
\[
\beta_1 \log n \leq c(u,v) \leq \beta_2 \log n. \qedhere
\]
\end{proof}



These results allow us to estimate how far the union of a minimum cut and a shortest path is from an optimal solution.
The diameter bound of Theorem~\ref{thm:random-diameter} and the connectivity bound of Lemma~\ref{lem:connectivity-bounds} show that the approximation ratio stays below $1+\epsilon$ on almost all inputs, so this simple algorithm is an average $(1+\epsilon)$-approximation scheme.

\begin{theorem}
\label{thm:approximation-scheme}
Let $G(n, p)$ be an Erdős–Rényi random graph with $p \geq \frac{\alpha \log n}{n}$, where $\alpha > 1$.
Then, for every $\epsilon > 0$, the algorithm that returns the union of a minimum $u$--$v$ cut and a shortest $u$--$v$ path is an average $(1+\epsilon)$-approximation scheme for the \MinCutPath{} problem.
\end{theorem}

\begin{proof}
The algorithm runs in polynomial time, since both a minimum cut and a shortest path can be computed in polynomial time~\cite{Cormen2022}.
By Lemma~\ref{lem:basic-bounds}, its output $S$ is a cut-path with $|S| \leq c(u,v) + d(u,v) - 1$.

Adding edges to a graph can only decrease the distances and can only increase the values $c(u,v)$.
Hence, by the standard coupling of the random graphs $G(n,p)$ for different edge probabilities, the bounds of Theorem~\ref{thm:random-diameter} and Lemma~\ref{lem:connectivity-bounds}, which are stated for $p = \frac{\alpha \log n}{n}$, remain valid for every $p \geq \frac{\alpha \log n}{n}$.
Thus, with probability tending to one, all pairs of distinct vertices $u,v$ satisfy
\[
d(u,v) \leq \diam(G) \leq \frac{\log n}{\log \log n}
\qquad\text{and}\qquad
c(u,v) \geq \beta_1 \log n.
\]
By Lemma~\ref{lem:basic-bounds}, the minimum cut is no larger than the minimum cut-path:
\[
c(u,v) \leq \cp(u,v) = \OPT(x).
\]

Combining these inequalities, we obtain
\[
\frac{|S|}{\OPT(x)} \leq \frac{c(u,v)+d(u,v)}{c(u,v)} = 1 + \frac{d(u,v)}{c(u,v)} \leq
1+\frac{\frac{\log n}{\log \log n}}{\beta_1 \log n} = 1 + \frac{1}{\beta_1 \log \log n}.
\]
The right-hand side is smaller than $1+\epsilon$ for all sufficiently large $n$.
Thus, the algorithm achieves the approximation ratio $1+\epsilon$ on a set of inputs whose probability tends to one, i.e., it is an average $(1+\epsilon)$-approximation scheme.
\end{proof}

\section{Conclusion}
\label{sec:conclusion}

We introduced the \MinCutPath{} problem and proved that its decision version is NP-complete.
In graphs of diameter two and in graphs with minimum cut-value at most two, $\cp(u, v) = c(u, v) + d(u, v) - 1$, and the union of a minimum cut and a shortest path is a minimum cut-path.
Dense Erdős–Rényi random graphs have diameter two with probability tending to one, and in sparser ones the same union is within a factor of $1+\epsilon$ of the optimum on almost all inputs.

Several questions remain open.
The first is to find further graph classes in which the problem can be solved in polynomial time.
A natural candidate is the class of graphs in which every pair of vertices is at distance at most two or has a minimum cut of size at most two; it contains both classes treated in this paper.

The second question concerns \emph{planar graphs}.
In a planar graph, every minimal edge cut corresponds to a cycle in the dual graph and vice versa, and several problems on cuts and flows can be solved faster in planar graphs than in general ones~\cite{itai1979maximum,henzinger1997faster}.
Whether this duality helps to solve \MinCutPath{} is open.

Third, this work is purely theoretical.
Experiments on real-world networks and on random instances would show how algorithms for \MinCutPath{} behave on average and could guide the design of heuristics.

Finally, graphs with a fixed value of $\cp(u,v)$ deserve a closer study, which may clarify the boundary between polynomially solvable and NP-hard instances.


\section*{Acknowledgments}

The author thanks Prof.~RNDr.~Martin Loebl,~CSc., who proposed the \MinCutPath{} problem, for his guidance and for valuable discussions.


\bibliographystyle{cas-model2-names}

\bibliography{references}

\end{document}
